\abstractPT  % Do NOT modify this line

U!REKA Play4Change é uma plataforma educacional gamificada e multiplataforma concebida para promover o envolvimento sustentado com temas de sustentabilidade e educação digital, desenvolvida como projeto de final de licenciatura no ISEL. A plataforma resolve dois problemas estruturais comuns ao software educativo existente: o estrangulamento na autoria manual de conteúdos, em que o esforço humano cresce proporcionalmente com a escala da plataforma, e a fragmentação dos ecossistemas móveis, em que manter aplicações nativas separadas para Android e iOS duplica o custo de engenharia. Ambos os problemas são resolvidos simultaneamente pela mesma base arquitetural.

O cliente móvel é construído com KMP (Kotlin Multiplatform) e Compose Multiplatform, disponibilizando uma aplicação Android completamente funcional e um alvo iOS a partir de uma única base de código partilhada. Novas tarefas educativas são produzidas automaticamente por um processo em lote noturno que invoca a API do modelo de linguagem Mistral através da biblioteca LangChain4j. Executar a geração num horário fixo, em vez de a pedido do utilizador, elimina a latência de inferência do caminho de pedido e desacopla completamente a disponibilidade da plataforma da disponibilidade do serviço de inteligência artificial. A repetição de conteúdos ao longo dos dias é evitada através da desduplicação semântica: cada tarefa gerada é representada como um vetor numérico de elevada dimensão e comparada com todos os vetores previamente armazenados usando similaridade de cosseno dentro do PostgreSQL com a extensão pgvector, descartando tarefas cujo significado é demasiado semelhante ao de conteúdos existentes, independentemente das diferenças de formulação.

A identidade dos jogadores é gerida através de autenticação delegada por OAuth~2.0 combinada com uma arquitetura de sessão de dois tokens: um token de acesso de curta duração mantido apenas em memória e um token de atualização de longa duração rotacionado em cada utilização, permitindo a deteção automática de credenciais roubadas. Um motor de pontuação do lado do servidor atribui recompensas configuráveis por conclusões corretas e atempadas e penalizações por inatividade. O backend é executado em Spring Boot e implantado como uma aplicação multi-contentor via Docker Compose. O modelo de custo em lote escala com o número de domínios temáticos em vez do número de utilizadores, tornando a plataforma progressivamente mais eficiente em termos de custo à medida que a adoção cresce.
